\chapter{Loading and visualising data}
\label{ch:carga}

With the calibration fixed (chapter~\ref{ch:calibracion}), loading a spectrum is
immediate: Fitbauer folds it, builds the velocity axis and displays it
normalised. This chapter describes the formats, the \emph{folding}, the
normalisation and the visualisation options.

\section{Supported formats}

\begin{center}
\begin{tabular}{@{}lp{0.62\textwidth}@{}}
\toprule
\textbf{Extension} & \textbf{Content} \\
\midrule
\ruta{.ws5}, \ruta{.adt} & \textbf{Raw counts} per channel (typ.\ 512
channels). They carry no velocity axis: they are folded and calibrated with the
fixed \param{Vmax}. \\[4pt]
\ruta{.csv}, \ruta{.txt}, \ruta{.dat}, \ruta{.exp} & Spectrum \textbf{already
folded} in velocity space (velocity/counts columns). They carry their own axis:
they are not folded and \emph{override} the calibration (§\ref{sec:calib-fija}). \\
\bottomrule
\end{tabular}
\end{center}

\begin{nota}
The \ruta{.ws5}/\ruta{.adt} files are lists of integers (one count per
channel). The velocity files (\ruta{.csv}/\ruta{.txt}/…) are two columns
—velocity and counts/transmission— separated by spaces, tabs or commas;
Fitbauer tolerates headers and comment lines. If a velocity file fails to
load, check that it has exactly two numeric columns per line.
\end{nota}

\section{Loading a spectrum}

\menu{File > Load…} opens the file selector. When you choose a raw-counts
spectrum, Fitbauer:

\begin{enumerate}[nosep]
  \item Reads the counts per channel.
  \item Detects the \textbf{folding centre} (or reads it from a NORMOS
        \ruta{.RES} file if one sits next to the spectrum).
  \item \textbf{Folds} the two branches of the triangular spectrum and
        \textbf{trims} a few channels from each end (§\ref{sec:folding} explains why).
  \item \textbf{Normalises} to a baseline of $\approx 1$ and estimates the Poisson noise.
  \item Builds the velocity axis with the \param{Vmax} of the fixed
        calibration.
\end{enumerate}

\autofig[\textwidth]{fig_load_alphafe}{The same $\alpha$-Fe \textbf{unfolded}
(left, 512 channels) and \textbf{folded} (right, in velocity). Unfolded, the
transducer measures the sextet \emph{twice} —once in each branch of the sweep—
mirrored about the \textbf{folding centre} (red line). Folding sums those two
halves: it reduces the noise and produces the symmetric spectrum on the right.}

\begin{nota}
The status bar and the file label report the number of channels, the detected
folding centre and the normalisation factor. \menu{File > Open recent}
reopens the most recently used spectra.
\end{nota}

\section{Folding and the centre}
\label{sec:folding}

The Mössbauer transducer runs through a \textbf{triangular} velocity cycle: each
velocity is measured twice, once on the rising branch and once on the falling
one. \emph{Folding} \textbf{sums those two halves} to average them: it improves
the signal-to-noise ratio and produces a single symmetric spectrum. The key
point is the \textbf{centre}: the (possibly fractional) channel at which it is
folded.

\begin{itemize}[nosep]
  \item Fitbauer detects the centre automatically (or reads it from a \ruta{.RES}).
  \item The \param{Centre} control of the calibration panel adjusts it by hand:
        the spectrum is \textbf{re-folded live} when it is changed.
  \item \menu{Fit > Find centre} refines it by minimisation.
  \item The \param{Fit folding point during the fit} checkbox includes it as a
        \textbf{free parameter of the discrete fit}: at each iteration the
        spectrum is refolded at the proposed centre, renormalised and the
        statistical weights are recomputed; when the fit finishes, the fitted
        value is written back to the \param{Centre} control and the data stay
        folded with it. Useful when the residual shows antisymmetric pairs
        around the lines; start from a reasonable value and keep few free
        parameters.
\end{itemize}

\paragraph{How Fitbauer folds.} Three details separate a careful folding from
a sloppy one, and all three are visible in the recovered linewidth.

\begin{itemize}[nosep]
  \item \textbf{Cubic interpolation.} The symmetry point almost never falls on a
        whole channel, so the paired channel has to be interpolated. Fitbauer
        uses a cubic B-spline. On synthetic spectra with a fractional folding
        point, the bias in the recovered linewidth drops from $+9.5\,\%$ with
        linear interpolation to $+2.8\,\%$ with the cubic one. NORMOS truncates
        to a whole channel and adds without interpolating.
  \item \textbf{Two search cycles.} The centre is located, refined in a second
        pass and only then used to fold.
  \item \textbf{Edge channels are kept.} Earlier versions trimmed a fixed number
        of channels from each end. They are no longer discarded as a matter of
        course: only \emph{anomalous} extreme channels are detected and trimmed
        --- the dead or saturated ones that show up now and then in the first or
        last channel of the multichannel analyser.
\end{itemize}

\begin{truco}
Dead edge channels are not a rare curiosity. Across a bank of real NORMOS jobs,
42 fits improved by up to a factor of 65 once they were detected, with the
reduced $\chi^2$ of the worst case falling from 102 to 1.45. If a fit refuses to
converge to anything sensible, look at the first and last channels of the raw
spectrum before blaming the model.
\end{truco}

\paragraph{Geometry effect.} The count rate is modulated by the \emph{position}
of the transducer, not only by its velocity. Fitbauer measures this modulation
and reports it, exactly as NORMOS does under \emph{Geometry effect} in its
\ruta{.RES}, but it does \textbf{not} correct it: the effect is antisymmetric
about the folding point, so folding itself cancels it. A relative amplitude
above $1\,\%$ betrays a spectrometer geometry problem and is worth investigating
before accepting the fit.

\begin{aviso}
A poorly chosen centre produces double or asymmetric lines. If you observe
spurious splittings or a shifted "ghost" spectrum, check the folding centre
\textbf{first} before touching the model parameters.
\end{aviso}

\subsection{Waveform: triangular or sinusoidal}
\label{sec:drive-form}

Everything above assumes a \textbf{constant-acceleration} drive, whose velocity
waveform is \textbf{triangular}: the velocity rises and falls in straight ramps,
so that it advances \emph{linearly} with the channel and the two halves of the
sweep cover the same velocities. That is why it makes sense to fold and use a
linear axis.

Some spectrometers use a \textbf{sinusoidal} drive instead. Then the velocity of
each channel is not linear, but
\[
  v_i \;=\; v_{\max}\,\sin\!\Big(\tfrac{2\pi\,(i-c_0)}{N}\Big),
\]
where $N$ is the number of channels and $c_0$ the \textbf{phase} (the channel at
which $v=0$). Since the axis is sinusoidal, folding by symmetry \emph{no} longer
gives a linear axis, and forcing a linear one would shift the lines.

\autofig[\textwidth]{fig_sine_drive}{Sinusoidal drive. Left: the velocity of each
channel follows a \emph{sine}, not a ramp. Right: the resulting spectrum, with
each channel placed at its real velocity. It is not folded: it is fitted as is.}

The \menu{Waveform} control of the calibration panel selects the mode:

\begin{itemize}[nosep]
  \item \textbf{Triangular} (default): it is folded and the axis is linear, as
        described above.
  \item \textbf{Sinusoidal}: following the NORMOS convention (\emph{do not fold}
        and fit the complete spectrum, with both branches at once), Fitbauer
        assigns each channel its real sinusoidal velocity and fits
        \textbf{without folding}.
\end{itemize}

Regarding the \textbf{phase} $c_0$: just as in the triangular case one has to
locate the centre, in the sinusoidal case one has to locate the symmetry point
of the sweep. Fitbauer autodetects it and exposes it in the \param{Centre}
control (adjustable live). Since a wrongly set phase would shift all the
velocities, it is also advisable to \textbf{refine it in the fit}: when
\emph{Fit centre} is enabled, $c_0$ enters as a \textbf{free parameter} and the
optimiser tunes it (just as \emph{Fit Vmax} tunes the amplitude). This is what
in NORMOS corresponds to \code{TRIANG=.FALSE.} with \code{FOLD=.FALSE.} and
\code{SIMULT=.TRUE.}

\begin{nota}
In sinusoidal mode the velocity axis is not monotonic (the same velocity appears
at several phases of the cycle), so the data are drawn as \textbf{points} that
run from $-v_{\max}\!\dots\!+v_{\max}$ and back. The model and the fit handle it
without problems, because they evaluate the model \emph{point by point} (they do
not need an ordered axis). The relaxation component is not supported with a
sinusoidal drive in this version.
\end{nota}

\section{Normalisation and noise}

So that fits are comparable between spectra, the data are \textbf{normalised} by
dividing by a factor (the 90th percentile of the folded counts), so that the
baseline sits at $\approx 1$ and the absorptions read as fractions. That same
factor is used to estimate the uncertainty of each point from the Poisson noise
(proportional to $\sqrt{N}$), which then weights the fit (see the
\emph{likelihood}, §\ref{sec:opciones-ajuste}).

\begin{nota}
The \textbf{baseline} and the \textbf{slope} (\param{Slope} control) of the
calibration panel model a possible drift of the background. They can be fixed or
left free in the fit like any other parameter.
\end{nota}

\section{Visualisation}

The plot area shows:

\begin{itemize}[nosep]
  \item The \textbf{data} (points), the \textbf{baseline}, the total
        \textbf{model} and, where appropriate, the \textbf{subspectra} of each
        component (with their optional area fill).
  \item Below, the \textbf{difference plot} (residual) data$-$model.
  \item In distribution mode, additionally the \textbf{envelope} of the
        distribution and each sharp component separately, and a panel with
        $P(\BHF)$.
\end{itemize}

The presentation is controlled from the \menu{View} menu:

\begin{itemize}[nosep]
  \item \menu{View > Show difference} / \menu{Show legend} /
        \menu{Show component area}: enable or hide each element.
  \item \menu{View > Plot style}: \emph{Classic}, \emph{Modern},
        \emph{Publication} or \emph{Dark}.
  \item \menu{View > Visual theme} and \menu{Colour theme}: general appearance of
        the interface.
  \item \menu{View > Configure panel layout…}: rearranges and saves layout
        presets.
\end{itemize}

The \textbf{toolbar} below the plot (Matplotlib) lets you zoom, pan, return to
the initial view and \textbf{save the figure} as an image (PNG, PDF, SVG).

\guicap{gui_plot_area}{Plot area with data, model, subspectra and residual, with
its navigation toolbar.}

\section{Comparing spectra}
\label{sec:comparar}

\menu{File > Compare spectrum…} overlays one or several reference spectra (with
their own colour palette) on top of the current one, useful for identifying
phases or comparing thermal treatments. \menu{File > Clear comparison} removes
all overlays. The comparison is purely visual: it does not alter the data or the
fit.

\section{Saving, exporting and reusing}

\begin{center}
\begin{tabular}{@{}p{0.42\textwidth}p{0.5\textwidth}@{}}
\toprule
\textbf{Action} & \textbf{What it saves} \\
\midrule
\menu{File > Save session…} & The \textbf{complete} state in \ruta{.json}
(data, fixed calibration, model, options, distribution). Reopen it with
\menu{Load session…}. \\[4pt]
\menu{File > Save fit…} & The current fit (parameters, curve,
distribution) in text. \\[4pt]
\menu{File > Export Markdown/PDF report…} & Readable report with parameters,
statistics and, where appropriate, the distribution table. \\[4pt]
\menu{File > Export reduced report…} & Abbreviated version of the report. \\
\bottomrule
\end{tabular}
\end{center}

\begin{truco}
The \textbf{session} is the recommended route for reproducibility: save it when
you finish an analysis and you will be able to resume it exactly as it was
—including the fixed calibration— at another time or on another computer. If you
have the laboratory API, \menu{File > Web} also lets you download
spectra/calibrations and \menu{Upload session JSON to web}.
\end{truco}

\section{NORMOS interoperability: \texorpdfstring{\texttt{.JOB}}{.JOB} files}
\label{sec:normos-job}

NORMOS is the program behind a large part of the published Mössbauer
literature, and many laboratories keep years of work stored as \ruta{.JOB}
files. Fitbauer reads and writes that format. It does \textbf{not} run NORMOS
and does not ship it: it only speaks its text format, which is not proprietary.

\subsection*{Importing}

\menu{File > NORMOS (.JOB) > Import NORMOS job…} reads a job and rebuilds the
model in Fitbauer's panels. \textbf{The spectrum is loaded too.} A \ruta{.JOB}
names its files in the first four lines with no path, because NORMOS ran under
DOS with everything in one directory; Fitbauer looks for the spectrum there,
ignoring upper/lower case, falling back to the job's own name and, as a last
resort, to the only spectrum in the folder. Keep every file of a job in the same
folder and the import works in one go.

Two families of job are recognised automatically:

\begin{itemize}[nosep]
  \item \textbf{NORMOS-SITE} (discrete sites). Each subspectrum becomes a
        component --- singlet, doublet or sextet --- with its values, its
        \param{fixed} boxes and its \code{NDEX}/\code{FACTOR}/\code{CONST}
        constraints.
  \item \textbf{NORMOS-DIST} (distributions). These open the
        $P(\BHF)$/$P(\dEQ)$ panel instead of the discrete mode. Their
        \code{NSUB} are \emph{not} sites: they are the points of a \textbf{grid}.
        Fitbauer translates the grid (origin and step), the shape (histogram,
        Gaussian, binomial or fixed), the $\IS(x)$ correlation and the edge
        anchors; the \emph{crystalline} subspectra become sharp components.
\end{itemize}

\subsection*{Exporting}

\menu{File > NORMOS (.JOB) > Export NORMOS job…} writes the current model in
NORMOS format. NORMOS has been verified to accept the file Fitbauer produces,
reproducing the original theory with a difference of exactly zero.

\subsection*{Convention conversions}

This is the delicate part, and getting it wrong is silent:

\begin{center}
\begin{tabular}{@{}p{0.26\textwidth}p{0.66\textwidth}@{}}
\toprule
\textbf{NORMOS} & \textbf{Meaning in Fitbauer} \\
\midrule
\code{WID}, \code{W13}, \code{W23} & \code{WID} is the width of lines 3,4 and
\code{W13}/\code{W23} are relative to it; Fitbauer's $\Gamma_1$ is the width of
lines 1,6. The conversion is $\Gamma_1 = \code{WID}\cdot\code{W13}$. \\[4pt]
\code{D13}, \code{D23} & \textbf{Area} ratios. Fitbauer's \param{int1}/\param{int2}
are \textbf{depth} ratios. They agree only when the widths are equal. \\[4pt]
\code{DEP} (or \code{ARE}) & The subspectrum \textbf{area} in mm/s, not a
depth. \\[4pt]
\code{NDEX}/\code{FACTOR}/\code{CONST} & Constraints in NORMOS's \textbf{global}
numbering, $13 + 15\,(n-1)$. \\
\bottomrule
\end{tabular}
\end{center}

\paragraph{The $\BHF$ scale.} NORMOS derives the sextet line positions from the
nuclear moments; Fitbauer uses the published $\alpha$-Fe pattern. They do not
differ by a simple scale factor. To reproduce a NORMOS $\BHF$ exactly, fit with
the NORMOS convention active (\code{core.constants.sextet\_pattern("normos")});
the difference is about $0.1$\,T.

\subsection*{The folding point is not imposed}

The \code{PFP} carried in a \ruta{.JOB} is the \textbf{seed} of the folding-point
search, not its result: NORMOS refines it over two cycles, and in real jobs it
ends up more than one channel away from what the file asked for. Fitbauer runs
its own search --- the correct counterpart --- and reports the \code{PFP} as
information only. Type it into \param{Centre} by hand if you want to force it.

There is a second subtlety. The refined point NORMOS prints in its \ruta{.RES}
is not where it folds either: its final routine truncates that value and adds
whole channels, placing the symmetry axis at $\lfloor \code{PFP} \rfloor + 0.5$.
Taking this into account is what makes its fits reproducible.

\subsection*{What is not carried over}

The importer warns about each of these, because what was \emph{not} translated
matters as much as what was:

\begin{itemize}[nosep]
  \item Czjzek / Le Caër distributions (\code{DISTRI=4}) and the
        Billard--Chamberod neighbour model (\code{METHOD=3}).
  \item Several overlapping distribution blocks: Fitbauer handles one.
  \item The \code{LAMDA} smoothing parameter. NORMOS's is absolute and
        Fitbauer's $\alpha$ is dimensionless, so there is no one-to-one
        conversion --- set it with the L-curve (section~\ref{sec:lcurve}). The \code{BETA}/\code{LAMDA} \emph{ratio} is preserved and
        becomes the edge anchor.
  \item \code{DTQ} in field distributions. NORMOS ignores it there, so carrying
        it over would introduce a correlation it never applied.
\end{itemize}

\begin{aviso}
Beware of inherited job files. The DIST format does not accept SITE keys such as
\code{NLINE}, \code{DEP}, \code{W13} or \code{W23}. If they were copied over from
another job, NORMOS reads and discards them \textbf{without a word}, so that
subspectrum never entered its fit at all. Fitbauer does warn about it.
\end{aviso}
